\documentclass[11pt]{article}
\usepackage[a4paper,margin=1in]{geometry}
\usepackage[T1]{fontenc}
\usepackage{lmodern,microtype}
\usepackage{amsmath,amssymb,amsthm}
\usepackage{booktabs,xcolor}
\usepackage[numbers,sort&compress]{natbib}
\usepackage[colorlinks=true,linkcolor=blue!55!black,citecolor=blue!55!black]{hyperref}
\newtheorem{theorem}{Theorem}[section]
\newtheorem{proposition}[theorem]{Proposition}
\newtheorem{remark}[theorem]{Remark}
\title{A Certified First-Omitted-Order Ladder for the Deterministic Tail}
\author{Bin Wang}\date{July 2026}
\begin{document}
\maketitle

\begin{abstract}
We organize the direct factorwise target-tail estimate into a certified ladder
of first omitted orders.  Arb replay gives explicit logarithmic and
multiplicative budgets for $N=13,21,29,37,45,53,61$.  Only $N=29$ matches the
currently archived order-28 deterministic head; higher rows are conditional
interfaces, not newly constructed heads.  The ladder is a target-side result
and leaves the cloud bridge and uniform quotient theorem open.
\end{abstract}

\section{Definition of the ladder}
Let $E_N$ denote the direct factorwise upper bound for
$\sum_{n\ge N}|a_n|/n$ from RH-264 \citep{WangRH264}.  We restrict to odd
$N$ so that the even trace threshold is $(N+1)/2$ and the odd endpoint series
starts at $N$.

\begin{proposition}[Monotone ladder]
For the fixed unit disk and the RH-13 certificate, the direct majorants satisfy
$E_{N'}<E_N$ whenever $N'>N$ are successive listed orders.  The same holds
for $\exp(E_N)-1$.
\end{proposition}

\begin{proof}
Each component is a positive tail of an absolutely convergent majorant.  The
first omitted order increases, so every component decreases; the exponential
map is increasing on the nonnegative reals.
\end{proof}

\section{Certified values}
The 100- and 200-decimal Arb replays evaluate the same outward interval
expressions independently.  The resulting safe endpoints are
\begin{center}
\small
\begin{tabular}{rcc}
\toprule
$N$ & $E_N$ & $\exp(E_N)-1$\\
\midrule
13 & $<2.3243607\times10^{-2}$ & $<2.3515845\times10^{-2}$\\
21 & $<8.73261\times10^{-4}$ & $<8.73642\times10^{-4}$\\
29 & $<2.6624745\times10^{-5}$ & $<2.6625100\times10^{-5}$\\
37 & $<9.32147\times10^{-7}$ & $<9.32147\times10^{-7}$\\
45 & $<4.432615\times10^{-8}$ & $<4.432615\times10^{-8}$\\
53 & $<1.705725\times10^{-9}$ & $<1.705725\times10^{-9}$\\
61 & $<7.175542\times10^{-11}$ & $<7.175542\times10^{-11}$\\
\bottomrule
\end{tabular}
\end{center}

\begin{theorem}[Current operational row]
The RH-253 head through order 28 and the $N=29$ row form a certified
deterministic target-side interface with logarithmic tail below
$2.6624745\times10^{-5}$.  No head through order 36 is present in the
archive, so the rows $N\ge37$ are conditional interfaces only.
\end{theorem}

\begin{proof}
The order-28 extent is the finite dictionary of RH-253 \citep{WangRH253};
the $N=29$ bound is RH-264.  The remaining rows are evaluated by the same
analytic majorant but require coefficients not claimed here.
\end{proof}

This ladder is not a cloud envelope, a selector, or a uniform quotient
theorem.  The five obligation vector remains $(0,0,0,1,1)$; all Gates A--E,
Hilbert--Polya claims, zeta-divisor equality, and RH implications remain
false/open.

\bibliographystyle{plainnat}\bibliography{references}
\end{document}
